\subsection{The Morita Correspondence of Modules}

In this subsection, the letters A and B denote Morita equivalent k-algebras and P an invertible \((A,B)_{k}\) -bimodule. Choose a \((B,A)_{k}\) -bimodule Q inverse to P and isomorphisms

\[
\lambda : \mathrm{P} \otimes_ {\mathrm{B}} \mathrm{Q} \rightarrow {} _ {s} \mathrm{A} _ {d} \quad \text { and } \quad \theta : \mathrm{Q} \otimes_ {\mathrm{A}} \mathrm{P} \rightarrow {} _ {s} \mathrm{B} _ {d}.
\]

For any left B-module V, denote by \(\theta_{V}\) the homomorphism of B-modules \(\theta\otimes1_{V}:Q\otimes_{A}P\otimes_{B}V\to V\) ; it is an isomorphism since \(\theta\) is an isomorphism. Likewise, for every left A-module M, denote by \(\lambda_{M}\) the homomorphism of A-modules \(\lambda\otimes1_{M}:P\otimes_{B}Q\otimes_{A}M\to M\) ; it is an isomorphism since \(\lambda\) is an isomorphism.

THEOREM 2 (Morita). --- a) Let V and W be left B-modules. The mapping \(g \mapsto 1_{\mathrm{P}} \otimes g\) is a bijection from \(\operatorname{Hom}_{\mathrm{B}}(\mathrm{V}, \mathrm{W})\) to \(\operatorname{Hom}_{\mathrm{A}}(\mathrm{P} \otimes_{\mathrm{B}} \mathrm{V}, \mathrm{P} \otimes_{\mathrm{B}} \mathrm{W})\) . The inverse bijection sends an element h of \(\operatorname{Hom}_{\mathrm{A}}(\mathrm{P} \otimes_{\mathrm{B}} \mathrm{V}, \mathrm{P} \otimes_{\mathrm{B}} \mathrm{W})\) to the element \(\theta_{\mathrm{W}} \circ (1_{\mathrm{Q}} \otimes h) \circ \theta_{\mathrm{V}}^{-1}\) of \(\operatorname{Hom}_{\mathrm{B}}(\mathrm{V}, \mathrm{W})\) .

\begin{enumerate}
\def\labelenumi{\alph{enumi})}
\setcounter{enumi}{1}
\tightlist
\item
  Every left A-module M is isomorphic to a module of the form \(\mathrm{P} \otimes_{\mathrm{B}} \mathrm{V}\) , where V is a left B-module.
\end{enumerate}

Let V and W be left B-modules. By Lemma 1 of VIII, p.~101, the mapping \(\varphi : g \mapsto 1_{\mathrm{P}} \otimes g\) from \(\mathrm{Hom}_{\mathrm{B}}(\mathrm{V}, \mathrm{W})\) to \(\mathrm{Hom}_{\mathrm{A}}(\mathrm{P} \otimes_{\mathrm{B}} \mathrm{V}, \mathrm{P} \otimes_{\mathrm{B}} \mathrm{W})\) is injective. By interchanging the roles of P and Q (and of A and B), we see that the mapping \(\psi : h \mapsto 1_{\mathrm{Q}} \otimes h\) from \(\mathrm{Hom}_{\mathrm{A}}(\mathrm{P} \otimes_{\mathrm{B}} \mathrm{V}, \mathrm{P} \otimes_{\mathrm{B}} \mathrm{W})\) to \(\mathrm{Hom}_{\mathrm{A}}(\mathrm{Q} \otimes_{\mathrm{A}} \mathrm{P} \otimes_{\mathrm{B}} \mathrm{V}, \mathrm{Q} \otimes_{\mathrm{A}} \mathrm{P} \otimes_{\mathrm{B}} \mathrm{W})\) is also injective. Now, the composition \(\psi \circ \varphi\) is the mapping \(g \mapsto \theta_{\mathrm{W}}^{-1} \circ g \circ \theta_{\mathrm{V}}\) . It is bijective, hence so is \(\psi\) . Consequently, \(\varphi\) is bijective, and its inverse is the mapping \(h \mapsto \theta_{\mathrm{W}} \circ (1_{\mathrm{Q}} \otimes h) \circ \theta_{\mathrm{V}}^{-1}\) .

Assertion b) follows from the fact that \(\lambda_{\mathrm{M}}\) is an isomorphism from \(\mathrm{P} \otimes_{\mathrm{B}} \otimes_{\mathrm{A}} \otimes_{\mathrm{M}}\) to \(\mathrm{M}\) .

Let V be a left B-module and W a submodule of V. Since the B-module P is projective (VIII, p.~101, Theorem 1), the canonical mapping from \(P \otimes_{B} W\) to \(P \otimes_{B} V\) is injective. We identify \(P \otimes_{B} W\) with its image in \(P \otimes_{B} V\) through this mapping. We adopt the analogous conventions when P and B are replaced by Q and A.

PROPOSITION 5. --- Let V be a left B-module. The mapping \(W \mapsto P \otimes_{B} W\) is an isomorphism from the set of B-submodules of V, ordered by inclusion, to the set of A-submodules of \(P \otimes_{B} V\) , ordered by inclusion. The inverse isomorphism sends an A-submodule N of \(P \otimes_{B} V\) to the image under \(\theta_{V}\) of the B-submodule \(Q \otimes_{A} N\) of \(Q \otimes_{A} P \otimes_{B} V\) .

Denote by \(D_{\mathrm{B}}(\mathrm{V})\) the set of B-submodules of V, ordered by inclusion, and define the sets \(D_{\mathrm{A}}(\mathrm{P}\otimes_{\mathrm{B}}\mathrm{V})\) and \(D_{\mathrm{B}}(\mathrm{Q}\otimes_{\mathrm{A}}\mathrm{P}\otimes_{\mathrm{B}}\mathrm{V})\) likewise. Let \(\varphi:D_{\mathrm{B}}(\mathrm{V})\to D_{\mathrm{A}}(\mathrm{P}\otimes_{\mathrm{B}}\mathrm{V})\) be the mapping \(W\mapsto P\otimes_{B}W\) and \(\psi\) the mapping from \(D_{\mathrm{A}}(\mathrm{P}\otimes_{\mathrm{B}}\mathrm{V})\) to \(D_{\mathrm{B}}(\mathrm{Q}\otimes_{\mathrm{A}}\mathrm{P}\otimes_{\mathrm{B}}\mathrm{V})\) given by \(N\mapsto Q\otimes_{A}N\) . These are increasing mappings, and the composition \(\psi\circ\varphi\) is the mapping \(W\mapsto\theta_{\mathrm{V}}^{-1}(\mathrm{W})\) , which is bijective. Consequently, \(\varphi\) is injective, and \(\psi\) is surjective. By replacing B with A and V with \(P\otimes_{B}V\) , we see that \(\psi\) is also injective. Hence, \(\varphi\) and \(\psi\) are bijective, and the inverse mapping of \(\varphi\) is indeed that described in the proposition.

Examples. --- 1) Let us apply Proposition 5 of VIII, p.~104, to the specific case \(\mathrm{V} = \mathrm{B}_s\) .

\begin{enumerate}
\def\labelenumi{\alph{enumi})}
\item
  The mapping \(J \mapsto PJ\) is an isomorphism from the ordered set \(\mathrm{D}(\mathrm{B}_{s})\) of left ideals of B to the ordered set \(\mathrm{D}(\mathrm{P})\) of A-submodules of P. The inverse mapping sends an A-submodule M of P to the left ideal \(\mathrm{J}(\mathrm{M})\) of B consisting of the elements b of B such that M contains Pb.
\item
  The mapping \(K \mapsto KP\) is an isomorphism from the ordered set \(\mathrm{D}(A_{d})\) of right ideals of A to the ordered set \(\mathrm{D}(P)\) of B-submodules of P. The inverse mapping sends a B-submodule V of P to the right ideal \(\mathrm{K}(V)\) of A consisting of the elements a of A such that V contains aP.
\end{enumerate}

Indeed, the A-module \(P \otimes_{B} B_{s}\) can be canonically identified with P. If J is a left ideal of B, then the canonical image of \(P \otimes_{B} J\) in \(P \otimes_{B} B_{s}\) corresponds to PJ through this identification. Consequently, the mapping \(J \mapsto PJ\) is an isomorphism of ordered sets from \(D(B_{s})\) to \(D(P)\) . Let \(J \in D(B_{s})\) . Denote the set of elements b of B such that PJ contains Pb by \(J'\) . It is a left ideal of B that contains J, and we have \(PJ' \subset PJ\) . Since the mapping \(J \mapsto PJ\) is an isomorphism of ordered sets, we necessarily have \(PJ' = PJ\) and \(J = J'\) . This proves a).

Assertion b) follows from assertion a) applied to the invertible \((\mathrm{B}^{\circ},\mathrm{A}^{\circ})_{k}\) -bimodule P.

Note that the ring A is a field if and only if the B-module P is simple.

\begin{enumerate}
\def\labelenumi{\arabic{enumi})}
\setcounter{enumi}{1}
\tightlist
\item
  Denote by \(\mathcal{B}_{\mathrm{A}}\) , \(\mathcal{B}_{\mathrm{B}}\) , and \(\mathcal{B}_{\mathrm{P}}\) the sets of two-sided ideals of A, two-sided ideals of B, and \((\mathrm{A},\mathrm{B})_k\) -sub-bimodules of P, respectively.
\end{enumerate}

\begin{enumerate}
\def\labelenumi{\alph{enumi})}
\item
  The mapping \(\mathfrak{b} \mapsto \mathrm{P}\mathfrak{b}\) is an isomorphism of ordered sets from \(\mathcal{B}_{\mathrm{B}}\) to \(\mathcal{B}_{\mathrm{P}}\) ; the inverse isomorphism sends an \((\mathrm{A},\mathrm{B})_k\) -sub-bimodule \(\mathrm{P}'\) of \(\mathrm{P}\) to the two-sided ideal of \(\mathrm{B}\) consisting of the elements \(b\) such that \(\mathrm{P}b \subset \mathrm{P}'\) .
\item
  The mapping \(a \mapsto aP\) is an isomorphism of ordered sets from \(B_{A}\) to \(B_{P}\) ; the inverse isomorphism sends an \((A, B)_{k}\) -sub-bimodule \(P'\) of P to the two-sided ideal of A consisting of the elements a such that \(aP \subset P'\) .
\end{enumerate}

Indeed, let J be a left ideal of B and \(P' = PJ\) . Then \(P'\) is an A-submodule of P and, by Example 1, the ideal J consists of the elements b of B such that \(Pb \subset P'\) . Moreover, \(P'\) is an \((A, B)_{k}\) -sub-bimodule of P if and only if J is a two-sided ideal of B. Thus a) follows from loc. cit.

Assertion b) follows from assertion a) applied to the invertible \((\mathrm{B}^{\circ}, \mathrm{A}^{\circ})_{k}\) -bimodule P.

PROPOSITION 6. --- Denote by \(\mathcal{B}_{\mathrm{A}}\) and \(\mathcal{B}_{\mathrm{B}}\) the sets of two-sided ideals of A and two-sided ideals of B.

\begin{enumerate}
\def\labelenumi{\alph{enumi})}
\item
  There exists an isomorphism of ordered sets \(f\) from \(\mathcal{B}_{\mathrm{A}}\) to \(\mathcal{B}_{\mathrm{B}}\) determined by the following property: if \(\mathfrak{a}\) is a two-sided ideal of \(\mathrm{A}\) and \(\mathfrak{b}\) a two-sided ideal of \(\mathrm{B}\) , then the relation \(f(\mathfrak{a}) = \mathfrak{b}\) is equivalent to \(\mathfrak{a}\mathrm{P} = \mathrm{P}\mathfrak{b}\) .
\item
  Suppose that the ring A is commutative, so that A can be identified with the center of B (VIII, p.~102, Corollary 1). The isomorphism \(f: \mathcal{B}_{\mathrm{A}} \to \mathcal{B}_{\mathrm{B}}\) sends an ideal \(\mathfrak{a}\) of A to the two-sided ideal \(\mathrm{Ba}\) of B, and we have \(\mathfrak{a} = \mathrm{A} \cap \mathrm{Ba}\) .
\end{enumerate}

Assertion a) follows from Example 2.

Now, suppose that A is commutative, and identify A with the center of B. Let a be an ideal of A. Then Ba is a two-sided ideal of B; we have PBa = aP and therefore \(f(a) = \text{Ba}\) . Let \(a'\) be the ideal \(A \cap Ba\) of A; it is contained in Ba and contains a, so \(Ba'\) is equal to Ba. Since f is bijective, it follows that \(a' = a\) .

Examples. --- 3) Let V be a left B-module. Then the correspondence given above sends the annihilator of the B-module V to the annihilator of the A-module \(P \otimes_{B} V\) . Indeed, denote the annihilator of the A-module \(P \otimes_{B} V\) by a and that of the B-module V by b. Let W be the \((A, B)_{k}\) -sub-bimodule of P consisting of the elements such that \(p \otimes v = 0\) for every element v of V. We have the inclusion \(Pb \subset W\) ; conversely, for every \(p \in W\) and \(q \in Q\) , we have that \(\theta(q \otimes p)\) belongs to b. Hence, the element b of \(D(B_s)\) corresponds to the element W of D(P). Likewise, \(a \in D(A_d)\) corresponds to W.

\begin{enumerate}
\def\labelenumi{\arabic{enumi})}
\setcounter{enumi}{3}
\tightlist
\item
  For any two-sided ideal \(\mathfrak{a}\) of A, denote the subset of \(\mathbf{M}_n(\mathrm{A})\) consisting of the matrices with entries in \(\mathfrak{a}\) by \(\mathbf{M}_n(\mathfrak{a})\) . It is a two-sided ideal of \(\mathbf{M}_n(\mathrm{A})\) . We have \(\mathbf{M}_n(\mathfrak{a})\mathrm{A}^n = \mathfrak{a}^n = \mathrm{A}^n\mathfrak{a}\) . It follows from Proposition 6 that every two-sided ideal of \(\mathbf{M}_n(\mathrm{A})\) is of the form \(\mathbf{M}_n(\mathfrak{a})\) , where \(\mathfrak{a}\) is a two-sided ideal of A.
\end{enumerate}

Remark. --- We keep the assumptions and notation above and suppose that the \((\mathrm{B},\mathrm{A})_{k}\) -bimodule Q is the dual \(P^{*}\) of the A-module P and that the isomorphisms \(\lambda\) and \(\theta\) are the canonical mappings \(\Lambda:P\otimes_{B}P^{*}\to_{s}A_{d}\) and \(\Theta:P^{*}\otimes_{A}P\to_{s}B_{d}\) (VIII, p.~101, Theorem 1). Since the A-module P is projective and finitely generated, we have a canonical isomorphism \(\vartheta_{M}:P^{*}\otimes_{A}M\to\operatorname{Hom}_{A}(P,M)\) for every A-module M (II, §4, No.~2, p.~271, Corollary). We leave it to the reader to reformulate the results of Subsections 3 and 4 by replacing the construction \(M\mapsto Q\otimes_{A}M\) with the construction \(M\mapsto\operatorname{Hom}_{A}(P,M)\) .

